\chapter{Pica}

\section*{Definition}
Persistent ingestion of non-nutritive, non-food substances for at least one month that is developmentally inappropriate and not part of a culturally sanctioned practice.

\section*{Pathophysiology}
The exact mechanism is unknown but is thought to involve iron or zinc deficiency dysregulating dopaminergic pathways in the striatum, leading to altered reward processing and craving for non-food items. Nutritional deficiencies may also impair oral mucosal enzymes that normally guide food selection.

\section*{Causes}
\begin{itemize}
  \item \textbf{Nutritional deficiencies:} Iron deficiency anemia (most common), zinc deficiency, calcium or magnesium deficiency
  \item \textbf{Developmental:} Autism spectrum disorder, intellectual disability, global developmental delay
  \item \textbf{Psychiatric:} Schizophrenia, obsessive-compulsive disorder, trichotillomania
  \item \textbf{Cultural/environmental:} Poverty, food insecurity, cultural practices (geophagy, amylophagy)
  \item \textbf{Medical:} Chronic renal disease (especially on dialysis), celiac disease, hookworm infection
\end{itemize}

\section*{When to See a Doctor}
See your doctor if:
\begin{itemize}
  \item Ingestion of toxic substances (lead paint, batteries, sharp objects)
  \item Evidence of heavy metal poisoning, gastrointestinal obstruction, or perforation
  \item Concurrent growth delay, weight loss, or severe anemia
\end{itemize}

\section*{Related Symptoms}
\begin{itemize}
  \item Fatigue, pallor, glossitis, koilonychia (spoon nails)
  \item Abdominal pain, constipation, vomiting
  \item Dental abrasion, halitosis
\end{itemize}
\textbf{Important:} Pica is associated with lead poisoning in children, especially when ingesting paint chips or soil in older housing.

\section*{Self-Care / Home Management}
\begin{itemize}
  \item Child-proof the home environment to restrict access to non-food items
  \item Provide appropriate oral stimulation (chew toys, crunchy foods) for sensory-seeking individuals
  \item Positive reinforcement for eating appropriate foods rather than punishing pica behavior
  \item Ensure balanced nutrition including iron-rich foods and consideration of supplementation
\end{itemize}

\section*{How Your Doctor Will Evaluate You}
\begin{itemize}
  \item Complete blood count with iron studies (ferritin, TIBC, transferrin saturation)
  \item Serum zinc, lead, and heavy metal levels
  \item Developmental and psychiatric assessment
  \item Abdominal imaging if obstruction or foreign body is suspected
\end{itemize}

\section*{Treatment Options}
\begin{itemize}
  \item Correct underlying nutritional deficiency; iron and zinc repletion often resolves pica
  \item Behavioral therapy including habit reversal training and differential reinforcement
  \item Treat comorbid psychiatric conditions with appropriate pharmacotherapy and therapy
  \item Monitor for complications: lead toxicity, bezoars, intestinal perforation, parasitosis
\end{itemize}

\clearpage
